\documentclass[10pt,a4paper]{amsart}
\usepackage[margin=19mm]{geometry}
\usepackage{amsmath,amssymb,mathtools}
\usepackage[scr=rsfs,cal=euler]{mathalfa}
\usepackage{xfrac}
\usepackage[unicode,hidelinks,bookmarks=false]{hyperref}
\newcommand{\ie}{i.e.\ }
\newcommand{\cf}{cf.\ }
\newcommand{\resp}{respectively\ }
\newcommand{\Subsection}{\subsection}
\providecommand{\nobreakdash}{\nobreak\hbox{-}}
\setlength{\emergencystretch}{2em}
\multlinegap0pt
\usepackage{bm}
\usepackage{bbm}
\usepackage{dsfont}
\usepackage{xspace}
\usepackage{cancel}
\usepackage{mathtools}
\usepackage{amsthm}
\makeatletter
\@ifundefined{theo}{\newtheorem{theo}{Theo}}{}
\@ifundefined{cor}{\newtheorem{cor}[theo]{Corollary}}{}
\@ifundefined{lemma}{\newtheorem{lemma}[theo]{Lemma}}{}
\@ifundefined{prop}{\newtheorem{prop}[theo]{Proposition}}{}
\makeatother
\newcommand{\Ga}{\mathbb{G}_a}
\newcommand{\Gm}{\mathbb{G}_m}
\newcommand{\G}{\ensuremath{\mathcal{G}}}
\newcommand{\ord}{\operatorname{ord}}
\newcommand{\ptweak}{{\rm{(P3)}}}
\newcommand{\pt}{{\rm{(P3$^u$)}}}
\newcommand{\sred}{{\sigma\text{-}\operatorname{red}}}
\newcommand{\s}{\ensuremath{\sigma}}
\title[Trivializing Torsors for Difference Algebraic Groups]{Trivializing Torsors for Difference Algebraic Groups}
\author{Annette Bachmayr, Michael Wibmer}
\date{Source: arXiv:2509.22095v1 (CC BY 4.0)}
\begin{document}
\maketitle

\noindent\emph{Source and licence.} Trivializing Torsors for Difference Algebraic Groups. Version arXiv:2509.22095v1 (CC BY 4.0), distributed under the Creative Commons Attribution 4.0 International licence, as recorded in the arXiv licence metadata for this exact version and shown on the abstract page https://arxiv.org/abs/2509.22095v1. Full rights notice: licenses/arxiv-2509.22095-CC-BY-4.0.txt.\par\smallskip

\noindent\textit{Editorial scope.} Counted unit: the Prop whose source label is \texttt{prop: sdim zero} in the article (source0.tex lines 2182--2184, source label \texttt{prop: sdim zero}), delivered together with the article's own complete original proof of it (source0.tex lines 2185--2218, one \texttt{proof} environment of 34 lines). No mathematics is written, repaired or re-derived: statement and proof are the article's own text line for line. Required context retained uncounted: prop:InductionPrinciple (lines 966--972); prop: order zero (lines 1344--1347); cor: reduce to connected (lines 1382--1383); prop: proof for diagonalizable (lines 1810--1812); lemma: reduce to sredued for strongly (lines 1851--1853); prop: vector groups of sdim 0 are strongly trivial (lines 1902--1905); prop: Zariski dense in simple (lines 2074--2076); lemma: embedd into Ga Gm or simple (lines 2109--2111). Every definition, display and label of those lines is retained unchanged. Omitted from this package and not redistributed: 1--965, 973--1343, 1348--1381, 1384--1809, 1813--1850, 1854--1901, 1906--2073, 2077--2108, 2112--2181, 2219--2565. Nothing inside a retained excerpt is omitted or altered: every retained body is delivered exactly as the article's lines are written, so no omission receipt of any kind accompanies this package. Citations: the retained excerpts carry no citation command at all, so no bibliography entry of the article is transcribed and no reference is redistributed. Reference adaptation: none was needed and none was made; every reference inside the delivered unit resolves inside it, so no unresolved reference or citation is produced. Printed slips kept literal: no printed word, symbol, hypothesis or number of the counted statement or of its proof was altered. Layout and class: the article's own class is \texttt{amsart} with packages including extdash, hyperref, geometry, xypic, enumerate, amsthm, amsmath, amssymb, amsfonts, todonotes; the wrapper is the lane's pinned \texttt{amsart} editorial wrapper at 10pt on A4, which restates the theorem-like environments the retained text needs and adds the article's own macro definitions that the retained text needs (\texttt{\textbackslash G}, \texttt{\textbackslash Ga}, \texttt{\textbackslash Gm}, \texttt{\textbackslash ord}, \texttt{\textbackslash pt}, \texttt{\textbackslash ptweak}, \texttt{\textbackslash s}, \texttt{\textbackslash sred}), with extra wrapper packages (bm, bbm, dsfont, xspace, cancel, mathtools). The delivered numbering is the unit's own. Assets: no retained span contains a figure environment, a graphics-inclusion command, a PDF-page inclusion, a table or a TikZ picture; no image file is copied into this package, the package declares no asset dependency and no image file is redistributed with it. Licence: version arXiv:2509.22095v1 of this article is distributed under the Creative Commons Attribution 4.0 International licence, as recorded in the arXiv licence metadata for this exact version and shown on the abstract page https://arxiv.org/abs/2509.22095v1. A full rights notice is delivered at licenses/arxiv-2509.22095-CC-BY-4.0.txt. Characters: every retained excerpt is pure ASCII, so the delivered engine, XeLaTeX with its default Latin Modern text font, reports no missing character.
\begin{prop}[Induction principle] \label{prop:InductionPrinciple}
Let $1\to N\to G\to H\to 1$ be an exact sequence of $\s$-algebraic groups over $k$.
\begin{enumerate}
    \item If $k$ is (uniformly) $H$-trivial and (uniformly) strongly $N$-trivial, then $k$ is (uniformly) $G$\=/trivial. 
    \item 	If $k$ is (uniformly) strongly $N$-trivial and (uniformly) strongly $H$-trivial, then $k$ is (uniformly) strongly $G$-trivial.
\end{enumerate}
\end{prop}

\begin{prop} \label{prop: order zero}
Let $G$ be a $\s$-algebraic group of order zero over an algebraically closed $\s$-field $k$. 
Then $k$ is uniformly strongly $G$-trivial. 
\end{prop}

\begin{cor} \label{cor: reduce to connected} Assume that $k$ is algebraically closed and let $G$ be a $\s$-algebraic group. If $k$ is (uniformly) strongly $G^o$-trivial, then $k$ is (uniformly) strongly $G$-trivial. 
\end{cor}

\begin{prop} \label{prop: proof for diagonalizable}
	Let $G$ be a diagonalizable $\s$-algebraic group of $\s$-dimension zero. If $k$ satisfies \ptweak{} (or \pt{}), then $k$ is (uniformly) strongly $G$-trivial.
\end{prop}	

\begin{lemma} \label{lemma: reduce to sredued for strongly}
	Let $k$ be an inversive $\s$-field and let $G$ be a $\s$-algebraic group over $k$. If $k$ is (uniformly) strongly $G_{\sred}$-trivial, then $k$ is (uniformly) strongly $G$-trivial.
\end{lemma}

	\begin{prop} \label{prop: vector groups of sdim 0 are strongly trivial}
 Let $G$ be a $\s$-algebraic vector group of $\s$-dimension zero. 
		If $k$ satisfies  \ptweak{} (or \pt{}), then $k$ is (uniformly) strongly $G$-trivial.
	\end{prop}

\begin{prop} \label{prop: Zariski dense in simple}
Let $G$ be a proper Zariski-dense $\s$-reduced $\s$-closed subgroup of a simple algebraic group $\G$. If $k$ satisfies \ptweak{} (or \pt{}) then $k$ is (uniformly) strongly $G$-trivial.
\end{prop}

\begin{lemma} \label{lemma: embedd into Ga Gm or simple}
	Assume that $k$ is algebraically closed. Let $G$ be a connected $\s$-algebraic group of positive order. Then there exists a proper normal $\s$-closed subgroup $N$ of $G$ such that $G/N$ is isomorphic to a Zariski-dense $\s$-closed subgroup of $\Ga$, $\Gm$ or a simple algebraic group.	
\end{lemma}

\begin{prop} \label{prop: sdim zero}
	Let $G$ be a $\s$-algebraic group of $\s$-dimension zero. If $k$ satisfies \ptweak{} (or \pt{}) then $k$ is (uniformly) strongly $G$-trivial.
\end{prop}

\begin{proof}
	We will prove this by induction on $\ord(G)$.
	The case when $G$ has order zero is treated in Proposition \ref{prop: order zero}. So we can assume $\ord(G)>0$.
	By Corollary \ref{cor: reduce to connected} we may assume that $G$ is connected. 
%	If $G$ is not almost-simple we can use Theorem \ref{theo: Jordan Hoelder for sreduced}, the induction hypothesis and Corollary \ref{cor: strong induction for tower} to conclude. 
%	Thus we can assume that $G$ is almost-simple. 
	According to Lemma \ref{lemma: embedd into Ga Gm or simple} there exists a proper normal $\s$-closed subgroup $N$ of $G$ such that $G/N$ is isomorphic to a Zariski-dense $\s$-closed subgroup of $\Ga$, $\Gm$ or a simple algebraic group. 
	As $\ord(N)<\ord(G)$, we can use the induction hypothesis together with Proposition \ref{prop:InductionPrinciple} (ii) to reduce to the case that $G$ is a Zariski-dense $\s$-closed subgroup of $\Ga$, $\Gm$ or a simple algebraic group. The case when $G$ is a $\s$-closed subgroup of $\Ga$ follows from Proposition \ref{prop: vector groups of sdim 0 are strongly trivial} and the case when $G$ is a $\s$-closed subgroup of $\Gm$ follows from Proposition \ref{prop: proof for diagonalizable}. So it only remains to treat the case when $G$ is a Zariski-dense $\s$-closed subgroup of a simple algebraic group. If $G_\sred$ is a proper $\s$-closed subgroup of $G$, then we can use the induction hypothesis and Lemma \ref{lemma: reduce to sredued for strongly} to conclude. Thus we can assume that $G=G_\sred$, i.e., $G$ is a proper $\s$-reduced $\s$-closed subgroup of a simple algebraic group. This case is treated in Proposition \ref{prop: Zariski dense in simple}.	
%	
%	
%	
%	
%	Because $G$ is almost-simple, $\ord(N)=0$. Thus Proposition \ref{prop: order zero} in combinations with
%	Proposition \ref{prop:InductionPrinciple}(ii) allows us to reduce to the case that $G$ is a Zariski-dense $\s$-closed subgroup of $\Ga$, $\Gm$ or a simple algebraic group.
%	
%	
%	
%	 Furthermore, by Lemma \ref{lemma: reduce to sredued for strongly} we may assume that $G$ is $\s$-reduced.
%%	
%	
%	Let $\G$ be an algebraic group such that $G$ embeds into $\G$ as a Zariski-dense $\s$-closed subgroup. Then $\G$ is connected ()
%	There exists a proper normal closed subgroup $\N$ of $\G$ such that $\G/\N$ is a simple algebraic group, $\Ga$ or $\Gm$.
%	
%	Now $N=\N\cap G$ is a normal $\s$-closed subgroup of $G$. Suppose $\N\cap G=G$, then $G\subseteq\N$. So $\N=\G$ because $\G$ is Zariski-dense in $\G$. Therefore $N$ is properly contained in $G$ and we have an embedding $G/N\to\G/\N$. 
%
%
%
%	
%	We would like to apply Proposition \ref{prop:InductionPrinciple}(ii) to the exact sequence $1\to N\to G\to G/N\to 1$.
%	Firstly, note that $(N^o)_\sred$ is normal in $G$ by Lemma \ref{lemma: connected and reduced remains normal}. As $(N^o)_\sred$ is connected and $\s$-reduced (Lemma \ref{lemma: Gsred is connnected}) we have $(N^o)_\sred=1$. In particular, $k$ is strongly $(N^o)_\sred$-trivial and it follows from 
%	Lemma \ref{lemma: reduce to sredued for strongly} and Corollary \ref{cor: reduce to connected} that $k$ is uniformly strongly $N$-trivial.
%
%	So it suffices to show that $k$ is (uniformly) strongly $G/N$-trivial. Now $G/N$ is a connected $\s$-reduced Zariski-dense $\s$-closed subgroup of either $\Ga$, $\Gm$ or a simple algebraic group. The $\Ga$-case is treated in Lemma \ref{lemma: vector groups of sdim 0 are strongly trivial}, the $\Gm$-case in \ref{lemma: diagonalizable groups of sdim 0 are strongly trivial} and finally the case of a simple algebraic group is treated in Proposition \ref{prop: Zariski dense in simple}. 	
\end{proof}

\end{document}
